\begin{abstract}
工具调用使智能体能够操作外部应用，也使用户信息随参数进入数据库、日志与第三方服务。现有评测主要考察工具选择、参数填写和任务完成情况，却较少区分"完成任务所必需的信息"与"工具虽然接受、但当前目的并不需要的信息"。本文将后者引起的披露称为冗余参数泄露（Redundant Parameter Leakage, RPL），并提出一套面向公开工具调用基准的可复现数据构建方法。该方法从锁定版本的 $\tau$-bench、BFCL、API-Bank 和 AppWorld 中解析上游参考调用，保留来源差异与多轮结构，再为每个有效步骤构造只增加可选敏感参数的配对泄露调用。最终数据包含 328 条任务，均通过结构校验、来源重放和执行检查，其中 321 条具有最终状态验证证据。本文以上游基准提供的参考调用作为功能参照，并采用统一的确定性规则生成参数权限标签，使构建结果可以复核和重复生成。该数据将冗余敏感信息披露与工具选择错误、必填参数缺失等一般调用错误区分开，为工具调用智能体的隐私评测提供了可执行、可追溯的实验样本。

\noindent\textbf{关键词：}工具调用智能体；隐私披露；数据最小化；冗余参数泄露；评测数据集；可复现研究
\end{abstract}
